% Minimal retained preamble for lane-p2 cards taken from
% sources/logic-fl05-ramsey-like-theorems-for-the-schre/pinned-original-source.tex
% (arXiv:2412.11598v3, CC BY 4.0).  This source ships no separate preamble file:
% the pinned tex is a concatenation of the authors' main.tex and macros.tex and
% carries its preamble inline, so the blocks below are sliced mechanically from
% that file and nothing has been renamed or redefined.
%
% Copied verbatim from the source:
%   line 2        : the authors' class line
%   lines 3-5,7,9 : graphicx, amsthm, the ams bundle, cleveref, url
%   line 10       : tikz-cd, used by the authors' commutative diagrams
%   lines 20-108  : the authors' macro definitions, taken in full so that every
%                   macro of the retained spans is available
%   lines 113-124 : the authors' theorem declarations
%   lines 126-128 : the authors' macros for the three calligraphic letters, which
%                   sit just after the theorem declarations in the source
%   lines 131-141 : the authors' machinery for restating a theorem
%   lines 143-148 : the authors' coloured comment macros
%
% Omitted source preamble lines and why:
%   line 1, 110-111 and 150 (concatenation comment, title, date, author): the
%       assembler supplies the title block, author list and date;
%   line 6 (hyperref) and line 8 (inputenc with the utf8x option): the forced
%       xelatex engine supplies the hyperref wiring itself and utf8x is
%       incompatible with it;
%   line 115 (numberwithin theorem by section): the assembler numbers the
%       retained material itself;
%   lines 120 and 121-123 are declarations that the authors left commented out or
%       that duplicate earlier ones and are not reproduced.
\documentclass{article}
\usepackage{graphicx}
\usepackage{amsthm}
\usepackage{amsfonts, amsmath, amssymb, amsthm}
\usepackage{cleveref}
\usepackage{url}
\usepackage{tikz-cd}
\newcommand{\Psf}{\mathsf{P}}
\newcommand{\Qsf}{\mathsf{Q}}
\newcommand{\Rsf}{\mathsf{R}}
\def\F{\mathcal{F}}
\def\U{\mathcal{U}}
\def\C{\mathcal{C}}
\def\I{\mathcal{I}}
\def\D{\mathcal{D}}
\def\L{\mathcal{L}}
\def\V{\mathcal{V}}
\def\A{\mathcal{A}}
\def\B{\mathcal{B}}
\def\P{\mathcal{P}}
\def\M{\mathcal{M}}
\def\Nc{\mathcal{N}}
\def\Q{\mathcal{Q}}
\newcommand{\cs}{2^{\NN}}
\newcommand{\bs}{\NN^{\NN}}
\newcommand{\bstr}{2^{<\NN}}
\newcommand{\uh}{\upharpoonright}
\newcommand{\converge}{\!\!\downarrow}
\newcommand{\diverge}{\!\!\uparrow}
\newcommand{\qvdash}{\operatorname{{?}{\vdash}}}
\newcommand{\nqvdash}{\operatorname{{?}{\nvdash}}}
\newcommand{\ISig}{\mathsf{I}\Sigma^0}
\newcommand{\Robinson}{\mathsf{Q}}
\newcommand{\IPi}{\mathsf{I}\Pi^0}
\newcommand{\IDelta}{\mathsf{I}\Delta^0}
\newcommand{\Iopen}{\mathsf{I_{open}}}
\newcommand{\BSig}{\mathsf{B}\Sigma^0}
\newcommand{\BPi}{\mathsf{B}\Pi^0}
\newcommand{\BDelta}{\mathsf{B}\Delta^0}
\newcommand{\BCSig}{\mathsf{BC}\Sigma^0}
\newcommand{\BCPi}{\mathsf{BC}\Pi^0}
\newcommand{\BCDelta}{\mathsf{BC}\Delta^0}
\newcommand{\LSig}{\mathsf{L}\Sigma^0}
\newcommand{\LPi}{\mathsf{L}\Pi^0}
\newcommand{\CSig}{\mathsf{C}\Sigma}
\newcommand{\RCA}[0]{\mathsf{RCA}}
\newcommand{\WKL}[0]{\mathsf{WKL}}
\newcommand{\WWKL}[0]{\mathsf{WWKL}}
\newcommand{\ACA}[0]{\mathsf{ACA}}
\newcommand{\ATR}[0]{\mathsf{ATR}}
\newcommand{\PIOOCA}[0]{\Pi^1_1\mathsf{-CA}}
\newcommand{\CAC}[0]{\mathsf{CAC}}
\newcommand{\ADS}[0]{\mathsf{ADS}}
\newcommand{\RRT}[0]{\mathsf{RRT}}
\newcommand{\FS}[0]{\mathsf{FS}}
\newcommand{\PFS}[0]{\mathsf{PFS}}
\newcommand{\TS}[0]{\mathsf{TS}}
\newcommand{\GRRT}[0]{\mathsf{GRRT}}
\newcommand{\LRRT}[0]{\mathsf{LRRT}}
\newcommand{\EM}[0]{\mathsf{EM}}
\newcommand{\SEM}[0]{\mathsf{SEM}}
\newcommand{\HalfEM}[0]{\mathsf{Half}\mbox{-}\mathsf{EM}}
\newcommand{\RT}[0]{\mathsf{RT}}
\newcommand{\SRT}[0]{\mathsf{SRT}}
\newcommand{\NN}[0]{\mathbb{N}}
\newcommand{\nat}[0]{\mathbb{N}}
\newcommand{\card}{\operatorname{card}}
\newcommand{\finsub}{\subseteq_{\mathtt{fin}}}
\newcommand{\W}{\mathrm{W}}
\newcommand{\sW}{\mathrm{sW}}
\newcommand{\lex}{\mathrm{lex}}
\newcommand{\base}{\operatorname{base}}
\newcommand{\range}{\operatorname{ran}}
\def\qt#1{``#1''}%

\newtheorem*{statement}{Statement}
\newtheorem{theorem}{Theorem}
\newtheorem{maintheorem}[theorem]{Main Theorem}
\newtheorem{lemma}[theorem]{Lemma}
\newtheorem{question}[theorem]{Question}
\newtheorem{proposition}[theorem]{Proposition}
%\newtheorem{statement}[theorem]{Statement}
\newtheorem{remark}[theorem]{Remark}
\newtheorem{definition}[theorem]{Definition}
\newtheorem{corollary}[theorem]{Corollary}
\newtheorem{example}[theorem]{Example}
\newcommand{\R}{\mathcal{R}}
\renewcommand{\S}{\mathcal{S}}
\newcommand{\T}{\mathcal{T}}

% Allow to repeat theorems
\makeatletter
\newtheorem*{rep@theorem}{\rep@title}
\newcommand{\newreptheorem}[2]{%
\newenvironment{rep#1}[1]{%
 \def\rep@title{#2 \ref{##1}}%
 \begin{rep@theorem}}%
 {\end{rep@theorem}}}
\makeatother

\newreptheorem{theorem}{Theorem}

% Comments
\usepackage{xcolor}
\newcommand{\ludovic}[1]{\textcolor{blue}{#1}}
\newcommand{\quentin}[1]{\textcolor{green}{#1}}
\newcommand{\lorenzo}[1]{\textcolor{orange}{#1}}
\newcommand{\oriola}[1]{\textcolor{red}{#1}}
\crefname{originalproposition}{original proposition}{original propositions}
\Crefname{originalproposition}{Original proposition}{Original propositions}
\makeatletter
\newlabel{rrt-omega-progressive-fs-omega}{{6.12}{source}{}{sourcepointer-1}{}}
\newlabel{rrt-omega-progressive-fs-omega@cref}{{[originalproposition][0][]6.12}{[1][0][]source}{}{}{}}
\makeatother
\crefname{originaltheorem}{original theorem}{original theorems}
\Crefname{originaltheorem}{Original theorem}{Original theorems}
\makeatletter
\newlabel{thm:pfs-omega-strong-cone-avoidance}{{6.14}{source}{}{sourcepointer-2}{}}
\newlabel{thm:pfs-omega-strong-cone-avoidance@cref}{{[originaltheorem][0][]6.14}{[1][0][]source}{}{}{}}
\makeatother
\title{Rainbow Ramsey for computably omega-bounded barriers admits strong cone avoidance}
\author{Lorenzo Carlucci, Oriola Gjetaj, Quentin Le Houérou, Ludovic Levy Patey}
\date{arxiv:2412.11598v3; CC BY 4.0}
\begin{document}
\maketitle
\noindent\textit{Editorial note.} The counted unit is the article's own unlabelled Corollary at source lines 1290-1293 together with its proof at source lines 1294-1302, a straight adjacent pairing in the pinned file. It is the corollary the article states right after its tightness proposition, and it transfers the avoidance theorem to rainbow colourings of barriers. The pinned file is the article's fixed version arXiv:2412.11598v3 under CC BY 4.0.\par\medskip
\noindent\textit{Editorial note.} This statement carries no source label, so the lane identifiers it by its opening line, exactly as it did for the unlabelled proposition of the same article: the claim key is the article identifier plus the statement's first line, and the statement span is the complete corollary environment.\par\medskip
\noindent\textit{Editorial note.} Because the statement has no label, the shared gate refuses the interval by design, so the lane ran its documented substitute checks instead: the live occupancy index reports the interval free, no card of any lane carries a token derived from the statement, and the public index carries exactly one entry for this source, at lines 585-587, which does not overlap this interval. The substitution and its outcome are recorded in the card's dedup block and in the lane STATUS.\par\medskip
\noindent\textit{Editorial note.} The retained proof names two results of the same article, the reduction from rainbow to progressive colourings at source lines 1033-1036 and the strong cone-avoidance theorem at source lines 1085-1088. Both are declared as non-load-bearing navigation pointers with the article's own section-wise numbers 6.12 and 6.14; the theorem is counted by another lane and only its statement is pointed at, with its proof left to the card that counts it.\par\medskip
\noindent\textit{Editorial note.} No citation occurs in the retained spans, so no bibliography entry is supplied. The retained preamble is the lane's shared slice for this source, 137 lines.\par\medskip
\noindent\textit{Allowed standard prerequisites (editorial).}\begin{itemize}
\item Barriers, rainbow colourings, progressive colourings and the cone-avoidance theorem. A barrier is a family of finite sets of natural numbers that is closed under taking initial segments and meets every infinite set exactly once along its increasing enumeration; a set is a rainbow for a colouring of a barrier when the sets it induces by initial segments take pairwise distinct colours, and the article's principle for barriers asks for such a rainbow. The same section proves that bounded progressive colourings of the barrier admit strong cone avoidance, and it proves separately that every colouring of a barrier can be turned computably into a constrained progressive colouring all of whose infinite free sets are rainbows for the original colouring. The author's own notions and those two results are the allowed inputs.. The counted corollary transfers strong cone avoidance from progressive colourings to rainbow colourings on computably bounded barriers, and it is quoted as the barrier-level form of the article's main avoidance theorem.
\end{itemize}\medskip
\ifdefined\hypertarget\hypertarget{sourcepointer-1}{}\fi
\noindent\textit{Source reference.} Original proposition 6.12: sources/logic-fl05-ramsey-like-theorems-for-the-schre/pinned-original-source.tex, source lines 1033--1036. Navigation pointer invoked by the retained proof for the reduction that turns a colouring of a barrier into a constrained progressive colouring whose free sets are rainbows. The cited result is an allowed input of the same section and its proof is not reproduced inside the span. The number is the article's own section-wise number for that proposition, counted over the numbered environments of section 6.\par\medskip
\ifdefined\hypertarget\hypertarget{sourcepointer-2}{}\fi
\noindent\textit{Source reference.} Original theorem 6.14: sources/logic-fl05-ramsey-like-theorems-for-the-schre/pinned-original-source.tex, source lines 1085--1088. Navigation pointer invoked by the retained proof for the strong cone-avoidance theorem for constrained progressive colourings, of which the counted corollary is the barrier-level consequence. That theorem is counted by another lane, so its proof is not reproduced here. The number is the article's own section-wise number for it, counted over the numbered environments of section 6.\par\medskip
\begin{corollary}
For every computably $\omega$-bounded barrier $B \subseteq [\NN]^{<\omega}$ and every~$k \in \NN^+$, $\RRT^B_k$ admits strong cone avoidance. 
%For every $k \in \NN^+$, $\RRT^{!\omega}_k$ admits strong cone avoidance.
\end{corollary}

\begin{proof}
Fix a set $Z$ and let $D$ be a non-$Z$-computable set.

Fix $k \in \NN^+$ and let $B \subseteq [\NN]^{<\omega}$ be a computably $\omega$-bounded barrier, with $h : \NN \to \NN$ the computable bound such that $B \subseteq [\NN]^{\leq h(\cdot)}$. Without loss of generality, we can assume that $h(x) \geq k$ for every $x \in \NN$. Let $f : B \to \NN$ be an instance of $\RRT^{B}_k$ and let $g : B \to [\NN]^{<\omega}$ be the $f'$-computable $k$-constrained (and hence $h$-constrained by our assumption that $g(x) \geq k$ for every $x \in \NN$) progressive coloring obtained in \Cref{rrt-omega-progressive-fs-omega}. 

$g$ can be extended to an $h$-constrained progressive coloring $\tilde{g} : [\NN]^{\leq h(\cdot)} \to [\NN]^{<\omega}$ be letting $\tilde{g}(s) = \emptyset$ for every $s \in  [\NN]^{\leq h(\cdot)} \setminus B$.

Then, by \Cref{thm:pfs-omega-strong-cone-avoidance}, there exists an infinite $\tilde{g}$-free set $G \subseteq \NN$ such that $D \not \leq_T G \oplus Z$. By construction of $g$, $G$ is also an infinite $f$-rainbow.
\end{proof}

\end{document}
